% ===== PRÉAMBULE CANONIQUE — à coller VERBATIM en tête de CHAQUE .tex =====
\documentclass[11pt,a4paper]{article}
\usepackage[utf8]{inputenc}\usepackage[T1]{fontenc}\usepackage{lmodern}
\usepackage[french]{babel}
\usepackage{amsmath,amssymb,mathtools}\usepackage{icomma}
\usepackage[version=4]{mhchem}          % \ce{...} : équations & formules
\usepackage{siunitx}\sisetup{locale=FR} % \qty{}{}, \si{}, \ang{}
\usepackage{chemfig}                    % \chemfig{...} : structures développées
\usepackage{xcolor}\usepackage{colortbl}
\usepackage{tikz}\usepackage{pgfplots}\pgfplotsset{compat=1.18}
\usepackage[most]{tcolorbox}\usepackage{enumitem}\usepackage{array,booktabs}
\usepackage[a4paper,margin=2.2cm]{geometry}
\usetikzlibrary{arrows.meta,calc,angles,quotes,positioning,patterns,backgrounds,shapes.geometric}
\definecolor{primary}{RGB}{222,49,99}\definecolor{primarydark}{RGB}{150,22,55}
\definecolor{defcol}{RGB}{0,114,178}\definecolor{methcol}{RGB}{52,92,148}
\definecolor{excol}{RGB}{0,135,90}\definecolor{attncol}{RGB}{200,40,55}
\tcbset{boxbase/.style={enhanced,breakable,boxrule=0.9pt,arc=2.5pt,
  left=3mm,right=3mm,top=2.3mm,bottom=2.3mm,fonttitle=\bfseries,coltitle=white}}
% --- boîtes TUTORIEL (noms EXACTS, ne pas renommer) ---
\newtcolorbox{cle}[1][]{boxbase,colback=defcol!6,colframe=defcol,
  title={\textsc{À retenir}\ifx\relax#1\relax\relax\else\textmd{ : \itshape #1}\fi}}
\newtcolorbox{methode}[1][]{boxbase,colback=methcol!7,colframe=methcol,sharp corners,
  title={\textsc{Méthode}\ifx\relax#1\relax\relax\else\textmd{ --- \itshape #1}\fi}}
\newtcolorbox{exemplebox}[1][]{boxbase,colback=excol!5,colframe=excol,
  title={\textsc{Exemple}\ifx\relax#1\relax\relax\else\textmd{ : \itshape #1}\fi}}
\newtcolorbox{piege}[1][]{boxbase,colback=attncol!6,colframe=attncol,
  title={\textsc{Piège}\ifx\relax#1\relax\relax\else\textmd{ : \itshape #1}\fi}}
\newtcolorbox{remarque}[1][]{boxbase,colback=black!4,colframe=black!45,
  title={\textsc{Remarque}\ifx\relax#1\relax\relax\else\textmd{ : \itshape #1}\fi}}
% --- boîtes EXERCICE / CORRIGÉ (noms EXACTS, auto-comptées) ---
\newtcolorbox[auto counter]{exo}[1][]{boxbase,colback=primary!4,colframe=primary,
  title={\textsc{Exercice}~\thetcbcounter\ifx\relax#1\relax\relax\else\textmd{ --- \itshape #1}\fi}}
\newtcolorbox[auto counter]{corrige}[1][]{boxbase,colback=excol!4,colframe=excol,
  title={\textsc{Corrigé}~\thetcbcounter\ifx\relax#1\relax\relax\else\textmd{ --- \itshape #1}\fi}}
\newcommand{\R}{\mathbb{R}}\newcommand{\N}{\mathbb{N}}\newcommand{\Z}{\mathbb{Z}}\newcommand{\ds}{\displaystyle}
% (un \usepackage{fancyhdr}+en-tête est possible ; il est ignoré côté web)
% =========================================================================
\begin{document}
\sisetup{per-mode=symbol}
\begin{center}\color{primarydark}\large\bfseries Thème 5 — Piles et électrolyse \quad$\bullet$\quad Niveau Facile\end{center}

\begin{exo}[anode, cathode et pôles]
Une pile associe les couples \ce{Zn^{2+}}/\ce{Zn} ($E_{\mathrm{r}}^{\circ}=-0{,}76$~V) et
\ce{Cu^{2+}}/\ce{Cu} ($E_{\mathrm{r}}^{\circ}=+0{,}34$~V).
\begin{enumerate}[label=\alph*)]
  \item Quelle électrode est l'anode ? la cathode ?
  \item Quel est le signe ($+$ ou $-$) de chacune ?
\end{enumerate}
\end{exo}

\begin{exo}[sens des électrons et du courant]
Dans la pile précédente (\ce{Zn} et \ce{Cu}), en fonctionnement :
\begin{enumerate}[label=\alph*)]
  \item Dans quel sens circulent les électrons dans le fil (de quelle électrode vers quelle autre) ?
  \item Dans quel sens circule le courant conventionnel $i$ ?
\end{enumerate}
\end{exo}

\begin{exo}[notation conventionnelle]
On considère la même pile zinc--cuivre.
\begin{enumerate}[label=\alph*)]
  \item Écris sa \textbf{notation conventionnelle} (anode à gauche).
  \item Que représentent le trait simple $|$ et le double trait $\big\|$ ?
\end{enumerate}
\end{exo}

\begin{exo}[charge et quantité d'électrons]
Un courant $I=\qty{1.5}{A}$ traverse une cuve pendant $t=\qty{20}{min}$. On donne
$F=\qty{96485}{C\per mol}$.
\begin{enumerate}[label=\alph*)]
  \item Convertis $t$ en secondes et calcule la charge $Q=I\,t$.
  \item Combien de moles d'électrons ont circulé ($n_{\ce{e^-}}=Q/F$) ?
\end{enumerate}
\end{exo}

\begin{exo}[premier calcul de Faraday]
On dépose de l'argent par électrolyse : \ce{Ag+ + e^- -> Ag} ($n=1$, $M(\ce{Ag})=\qty{108}{g\per mol}$).
Le courant est $I=\qty{1.0}{A}$ pendant $t=\qty{965}{s}$. Calcule la masse d'argent déposée
$m=\dfrac{M\,I\,t}{n\,F}$ (avec $F=\qty{96485}{C\per mol}$).
\end{exo}

\end{document}
